% Folder 134-2, batch 04 — archivists' pages 61–80.
% Pass: Opus 5.5 (claude-opus-5-5), 2026-10-03 — first pass, unchecked against the pages by a human.
%
% Special rules for folder 134-2 (issue #44) applied: Pursuing Stacks text is not
% transcribed; Grothendieck's 1975 letters to Lawrence Breen are; third-party
% letters are not. No third-party letter in this batch.
%
% Letters in this batch:
%  - pp. 61–77 (letter pp. « -16- » to « -32- », renumbered by hand 46 to 61 and
%    « 61 bis »): continuation of « Villecun 17/19 July 1975 », to Lawrence Breen
%    (« Dear Larry »), English, typed with his handwritten corrections, marginal
%    « (App. 12) » to « (App. 18) » marks and footnote calls (28)–(32). The letter
%    ENDS on p. 77 (« Very cordially yours, Alexandre Grothendieck. »). Transcribed.
%    The letter's own pp. 2–11 (missing from batch 3) do not appear in this batch.
%  - pp. 78–80 (his pp. 62–64): typescript « NOTES to Chapter I and Appendix »,
%    notes (1) to (25), part of the Pursuing Stacks typescript (1983; note (11)
%    is « Added 23.2.83 »). Not transcribed: \page{} plus one \note{} each, per
%    the special rules. Note (25) runs on past p. 80.
\documentclass[11pt,a4paper]{article}
\input{../preamble/grothendieck}

\folder{134-2}
\batch{4}
\pages{61}{80}
\foldertitle{[Chapitre I : Take off (pages 1 à 65) et table des matières provisoire] : tapuscrits annotés (19/02-22/02/1983), note manuscrite (s.d.), copies de lettre (1975, s.d.).}
\dating{1975-[1983]}
\watermark{Édition de démonstration}

\begin{document}

\page{61}
\note{suite de la lettre à L. Breen, « Villecun 17/19 July 1975 », commencée avant ce lot ; la phrase de la p. 60 (« with the exception of ») se poursuit ici. Pagination de la lettre : « -16- », surchargée en 46.}
$J_0$ if I remember well, ($J_0$ had as abelian part the abelian part of $\mathrm{Alb}_{X/k}$ the \struck{ordinary} \add{usual} generalised jacobian). It\add{'}s construction, inspired by the residual complex, passes by generalised jacobians (in an appropriate cohomological sense) of the \struck{localised} \add{localizations} $\operatorname{Spec} \underline{O}_{X,x}$ of $X$ at its different point.
\note{« Alb » : la frappe porte « A », la fin du mot est tracée à la main.}

N.B. $\underline{H}_0(J_*)$ was the « generalised Jacobian » of $X$, i.e. there existed a homeomorphism $X \to \underline{H}_0$, which was universal for homomorphisms of $X$ into \add{comm.} locally proalgebraic groups. For $X$ connected, $\underline{H}_0$ is an extension of $\mathbb{Z}$ by an appropriate proalgebraic group.
\note{l'alinéa « N.B. … » est encadré à la main et une flèche le renvoie après la phrase suivante (« It is possible … smooth. ») : il semble qu'il doive être déplacé. « homeomorphism » est la lecture de la frappe.}

It is possible that, at first, I restricted to the case of $X$ smooth. The cohomology role of this complex was that of a complex of \emph{homology}
\marginal{formule sur une ligne}
\[
(*) \qquad H^i(X, G_X) \simeq \mathbb{E}\mathrm{xt}^i(J_{*X/k}, G)
\]
but for which coefficients \add{$G$}? I believe I took arbitrary commutative algebraic groups $G$ but worked with the Zariski topology (malédiction !)\add{.} Even in the case of discrete $G$, I considered the Zariskian $H^i$, this gives slightly stupid cohomology groups, evidently. I realised that one should work ultimately in étale cohomology, and that the construction of the $(J_{\cdot})_{X/k}$ will evidently be modified \add{accordingly}. As for the significance of the $\mathbb{E}\mathrm{xt}^i$ (hypercohomology), at a moment where Serre had developed the formalism for proalgebraic groups, one was not too fearful of taking it in the category of such objects -- \struck{the case of} \add{and \ill{} in the sense of a} « derived category » which at that moment had never yet been \add{explicitly} defined \struck{or} \add{and} studied\add{.} (We have, after all, somewhat progressed since those days!). I have the impression, in view of these antique cogitations, heuristic as they were, that it should now be possible to develop at present such a theory of $J_{*X/k}$, in cohomology f.p.p.f., giving a formula $(*)$ without limitation on the degree $i$ of the cohomology. (N.B. But $J_*$ evidently no longer stops in $\dim X = n$ but in $\dim 2n$. It is nevertheless possible that the components $J_i$ might be of $\dim 0$ for $i > n$).
\note{la note marginale, lue « formule sur une ligne » (l'ordre des mots tracés en arc est incertain), et les crochets qui isolent la formule $(*)$ sont des consignes de mise en page. Dans « $(J_{\cdot})_{X/k}$ » l'indice est surchargé et lu avec doute. Après « and » un mot biffé et noirci est illisible. Au-dessus de « f.p.p.f. » un trait de crayon effacé.}

I believe that the construction of the $J_*$ does not commute with base change, but merely does so in the derived category sense.
\marginal{laisser espace}

\page{62}
\note{pagination de la lettre : « -17- », surchargée en 47.}
\marginal{(App.~12)}
(D) Let $X/k$ be a smooth scheme (for simplicity) \struck{on} \add{over} a field $k$, separated and of finite type, of relative dimension $d$, and $n$ an integer $>0$. If $n$ is prime to the characteristic and if $F$ is a sheaf of coefficients on $X$ which is annihilated by $n$, the global duality tells us that $Rf_!(F)$ and $Rf_*(R\,\mathit{Hom}(F, \mu_n^{\otimes d}))$ ($\mu_n$ = sheaf of $n$th roots of unity $= \ker(\mathbb{G}_m \xrightarrow{n} \mathbb{G}_m)$) are dual to each other with values in $(\mathbb{Z}/n\mathbb{Z})_k$, for example $Rf_!(\mathbb{Z}/n\mathbb{Z})$ and $Rf_*(\mu_n^{\otimes d})$, or $Rf_!(\mu_m^{\otimes})$ and $Rf_*(\mathbb{Z}/n\mathbb{Z})$, are dual to each other -- at least with a shift of amplitude $2d$ in dimension. (As $\mathbb{Z}/n\mathbb{Z}$ is injective over itself, this gives in fact perfect dualities
\[
R^if_!(F) \times R^{2d-i}f_*(\mathbb{R}\,\mathrm{hom}(F, \mu_n^{\otimes d})) \longrightarrow \mathbb{Z}/n\mathbb{Z} \,.)
\]
\note{« $\mu_m^{\otimes}$ » : ainsi dans la frappe (indice $m$, exposant sans $d$).}

If now one no longer assumes $n$ prime to the characteristic, for example $n$ is a power of $p$ = characteristic of $k > 0$), it seems that everything collapses: to start with, one no longer knows (for $d > 1$) by what to replace $\mu_n^{\otimes d}$ …

The extraordinary miracle is that for $d = 1$, i.e. $X$ a smooth curve, everything continues to work perfectly, provided one states things with care! The first verifications are made for example with $F = \mathbb{Z}/p\mathbb{Z}$, $\mu_p$, \add{or} $\alpha_p$, \add{with $X$} complete -- one finds it's O.K. by virtue essentially of the autoduality of the jacobian. One can make these examples more sophisticated on taking \emph{twisted} coefficients, and $X$ not complete -- one convinces oneself this works always! Simply, it is necessary to note that here the $R^if_*(F)$, $R^if_!(F)$ have a « continuous » structure (they are essentially proalgebraic groups). This corresponds to the well known phenomenon in class field theory that the structure of $\pi_{1\mathrm{ab}}$ of $X$, when $X$ is not complete, is \emph{continuous} -- hence same holds for $H^1(X, \mathbb{Z}/p^n\mathbb{Z})$ \add{say}.
\note{dans « $\mathbb{Z}/p\mathbb{Z}$ » le $p$ est retracé à la main ; dans « $H^1(X, \mathbb{Z}/p^n\mathbb{Z})$ » la frappe portait « $\mathbb{Z}/n\mathbb{Z}$ », corrigé à la main en $p^n$, correction répétée en marge (« a/ $p^n\mathbb{Z}$ »). Le début de « same » est retouché.}

By the way, I point out for you that Serre once proposed (without ever writing it down, I think) a theory of duality for commutative unipotent algebraic groups, \emph{mod radical isogeny}, duality with values in $\mathbb{Q}/\mathbb{Z}$ (or $\mathbb{Q}_p/\mathbb{Z}_p$). He found that if (when $k$ is algebraically closed, say) $G$ is

\page{63}
\note{pagination de la lettre : « -18- », surchargée en 48.}
such a group, then $G' = \mathrm{Ext}^1(G, \mathbb{Q}/\mathbb{Z})$ can canonically be given a structure of quasi-algebraic group (i.e. defined mod radical isogeny), doubtless in a unique manner provided it verifies some functorial properties, and on requiring that for $G = \mathbb{G}_a$ one finds that $\mathrm{Ext}^1(\mathbb{G}_a, \mathbb{Q}/\mathbb{Z}) \simeq \mathbb{G}_a$ with the usual structure. Let $\Delta G = G' = \mathit{Ext}^1(G, \mathbb{Q}/\mathbb{Z})$. One finds $G \simeq \Delta\Delta G$ i.e. $\Delta$ is an authentic autoduality! I call $\Delta$ \emph{Serre duality}. It surely goes over to ind-progroups on an arbitrary base field (not necessarily algebraically closed) in the case $p > 0$. Moreover, for finite étale groups, it is $\mathit{Ext}^0(G, \mathbb{Q}/\mathbb{Z})$ (Pontrjagin duality) which gives a perfect duality. One could \struck{assemble} \add{screw together}, in an appropriate derived category, Serre duality and Pontrjagin duality, by taking $G \longmapsto \Delta G = \mathbb{R}\mathit{Hom}(G, \mathbb{Q}/\mathbb{Z})$: one calls this (« cohomological ») Serre duality. This will be a magnificent autoduality, if one puts oneself in a derived category where the $\underline{H}^i$ of the envisaged complexes are (up to passing to the limit) extensions of étale groups by connected unipotent groups. Now one \struck{only meets} \add{gets only} such complexes, \struck{on} \add{by} « integrating » finite coefficients $F$ on $X$ by $Rf_!$ or $Rf_*$. This being said, by \struck{taking} \add{passing} to the limit in the initial formulation (or equivalently by replacing the $(\mathbb{Z}/n\mathbb{Z})_k$, previously considered, by $(\mathbb{Q}/\mathbb{Z})_k$ on $k$, and forming $f^!(\mathbb{Q}/\mathbb{Z})_k = (\mu_\infty)_X$) the duality formula takes the form
\[
\Delta(Rf_!(F)) \simeq Rf_*(DF[2]) \qquad \text{« shift » of dimension}
\]
where $D$ is the « Cartier duality » $R\mathit{Hom}(F, \mu_\infty)$ (or $R\mathit{Hom}(F, \mathbb{G}_m)$ if one prefers?), and $\Delta$ is the Serre duality: cohomology with proper supports and \add{with} arbitrary supports are exchanged by duality, when one takes upstairs Cartier duality, and downstairs Serre duality.
\note{dans « $G' = \mathrm{Ext}^1$ » (deux fois) l'exposant frappé est biffé et remplacé à la main par un accent, répété en marge (« $G'$ »). Dans « $(\mu_\infty)_X$ » et « $\mu_\infty$ » l'indice frappé est noirci et « $\infty$ » écrit à la main au-dessous. Une flèche à la main relie « shift » au $[2]$ de la formule.}

The validity of the duality formula is not open to doubt -- the principal work for establishing it consists certainly in a careful description of \add{the} \struck{coefficient} categor\struck{y}\add{ies} \add{of coefficients} with which one is working, as well on $X$ as on $k$, and of the functors $D$ and $\Delta$. As the definition of an arrow is immediate, once the building of the machine has been accomplished, the validity of the

\page{64}
\note{pagination de la lettre : « -19- », surchargée en 49.}
formula should result without difficulty from the usual « dévissages » which allow one to verify the duality in the particular standard cases $F = \mathbb{Z}/p\mathbb{Z}$, $\mu_p$, $\alpha_p$ on a smooth, complete $X$. (N.B. the case of coefficients prime to the characteristic is already known.) Let us make explicit what the formula of duality says for $R^1f_*(G_X)$, where $G$ is a finite group étale on $k$ (the most important case being $G = (\mathbb{Z}/p^m\mathbb{Z})_k$); one recovers \struck{the} \add{Serre's} description \struck{of Serre} of « geometric classfield\add{ theory} », in \struck{the form of} \add{terms of} extensions by $G$ of a generalised jacobian of $X$. Thus, the duality formula can be understood as a cohomological version, considerably enriched, of geometric class field theory. When the base field $k$ is finite, to retrieve the classfield theory in the classical form, one can use « the trick of Lang » [on the relation between the « arithmetic » $\pi_1$ of a smooth, connected commutative algebraic group $J$ on $k$ and its $H^0(k, J) = J(k)$: the $\pi_1^{\mathrm{ar}}(J)$ classifies the isogenies above $J$ with kernel a constant group $\pi_1^{\mathrm{ar}}(J) \simeq H^0(k, J)$] -- in its cohomological form, which may be stated:
\[
\Delta_0 \mathbb{R}\Gamma_k(J^*) \simeq \mathbb{R}\Gamma_k(\Delta J^*[1]) \,,
\]
where $\Delta$ is Serre duality, $\Delta_0$ Pontrjagin duality for the totally disconnected topological abelian groups (duality with values in $\mathbb{Q}/\mathbb{Z}$), $J^*$ a complex of algebraic ind-progroups on $k$. Taking account of th\struck{e}\add{is} \struck{« formula of Lang duality »} \add{« Lang duality formula »} and applying $\mathbb{R}\Gamma_k$ to the formula of duality for geometric classfields, one \struck{finds} \add{gets the} « duality formula \struck{for} \add{of} arithmetic classfield\struck{s} \add{theory} »:
\[
\Delta_0(H_!(X, F)) \simeq H^*(X, D(F)[3])
\]
(isomorphism of totally disconnected topological groups).
\note{« $G_X$ » : la frappe porte le même $G$ évidé que dans $\mathbb{G}_a$ ; d'après la suite (« where $G$ is a finite group étale »), il s'agit du faisceau constant $G_X$.}

Another remark: when $F$ \struck{comes from an} \add{is not an} « étale sheaf », but has a continuous structure such as $\alpha_p$, one must be careful in the definition of $Rf_!(F)$, for $X$ non complete, starting from the compactification $\tilde{X}$; thus, if $F$ comes from an « admissible » sheaf $\tilde{F}$ on $\tilde{X}$, one must have an exact triangle

\page{65}
\note{pagination de la lettre : « -20- », surchargée en 50.}

\begin{tikzcd}[column sep=small]
 & Rf_!(\hat{\tilde{F}}) \arrow[dl] & \\
Rf_!(F) \arrow[rr] & & R\tilde{f}_*(\tilde{F}) \arrow[ul]
\end{tikzcd}

where $\hat{\tilde{F}}$ is the \emph{formal completion} of $\tilde{F}$ along $\tilde{X} - X$ (a finite number of points …). It is here, unless I am mistaken, that appears the link with local class field theory, in its cohomological version, on which I am going now to say \struck{several} \add{a few} words.
\marginal{espacer}

\marginal{(App.~13)}
(E) \struck{Local} \add{Local} class field\struck{s} \add{(theory)} as a duality formula
\note{l'intitulé est souligné dans la frappe ; il était précédé d'un appel frappé biffé (illisible) ; « (E) » et « Local » sont récrits à la main, « (theory) » ajouté.}

Let $V$ be a complete discrete valuation ring \add{with} residue field $k$ -- assume either that $k$ \struck{can be} \add{has been} lifted to $k \subset V$ (and therefore $V \simeq k[[T]]$) or that $k$ is perfect \add{of char.~$p>0$}. In order to fix ideas, and to be sure that I'm on solid ground, I consider at first on $K$ (= the field of fractions of $V$) \emph{finite} coefficients $F$ (as on $X$ previously) and I consider \add{the objects} $H^i(K, F)$, or $R\Gamma_K(F)$. The main work to be done consists in defining an adequate category of coefficients over $k$ (perhaps the same one as in (D)) and a functor
\[
F \longmapsto \mathbb{R}\underline{\Gamma}_K(F)
\]
with values in the category of such coefficients, in such a manner that the following isomorphism holds.
\[
R\Gamma_K(F) \simeq R\Gamma_k(R\underline{\Gamma}_K(F)) \,.
\]
\note{« $H^i(K, F)$ » : l'exposant, surchargé, est lu avec doute.}

This corresponds to the intuition (acquired directly from elementary examples) according to which for $k$ algebraically closed, say, the $H^0(K, F)$, $H^1(K, F)$ … are endowed with a structure of $k$-algebraic group (ind-pro …). In this construction, the ring scheme of Witt vectors over $k$ (introduced by Serre) and the « Greenberg functor » (associating to a $V$-scheme a $k$-prescheme) will play an essential role.
\note{« endowed » : la frappe est corrigée à la main, correction répétée en marge (« w/ »).}

This being done, the duality formula will be formally stated as in (D) above:
\[
\Delta R\underline{\Gamma}_K(F) \simeq R\underline{\Gamma}_K(DF[1])
\]
where $D$ stands for Cartier duality, $\Delta$ for Serre duality. When the

\page{66}
\note{pagination de la lettre : « -21- », surchargée en 51.}
residue field is finite, it becomes (via « Lang's trick » mentioned previously)
\[
\Delta_0 R\Gamma_k(F) \simeq R\Gamma_K(DF[2])
\]
$\Delta_0$ standing for Pontrjagin duality. The formula contains local geometric class field theory à la Serre, and arithmetical local class field theory in its classical form.
\note{dans la première formule de la page, la frappe porte $\Gamma_k$ à gauche et $\Gamma_K$ à droite ; laissé tel quel.}

\emph{Remarks}

(a) If $F$ is prime to the residue characteristic the formula is very easy to prove and well known. It may be considered a very special case of the « induction formula » for a morphism $i : s \longmapsto S$, in the duality formalism:
\[
i^!(D_S(F)) = D_S(i^*(F))
\]
(we take here the inclusion of $p = \operatorname{Spec}(k)$ in $S = \operatorname{Spec}(V)$). Thus the work to be done concerns the $p$-primary coefficients, for $p = \operatorname{char} k > 0$. The most subtle case is that of unequal characteristics.

(b) The functor $R\underline{\Gamma}$ may be obtained by composing $\mathbb{R}j_*$ (where $j : U = \operatorname{Spec}(K) \longrightarrow \operatorname{Spec}(V) = S$ is the inclusion) with a cohomological version of the « Greenberg functor ».

(c) In (D) and (E), I restricted myself to finite coefficients $F$ -- it's for those that I am sure of what I assert. But it is certainly true that the duality formula is even richer, that something may still be asserted for example for $F$ a not necessarily finite group scheme, for example an abelian scheme (with a few degenerate fibres in the case of (D)?), but I have never entirely clarified this question, even on a heuristic basis. I vaguely recall a formula which should be contained in the formalism (say if $k$ is algebraically closed): for $F$ an abelian scheme on $K$, $F'$ the dual abelian scheme and $G'$ the pro-algebraic group over $k$ attached « à la Greenberg » to its Neron model, then one has
\[
H^1(K, F) \overset{?}{\simeq} \mathrm{Ext}^1_{k\text{-grp}}(G', \mathbb{Q}/\mathbb{Z})
\]
(N.B. without any guarantee.) In principle, the previously mentioned duality conjecture concerning Neron models of SGA~\uncertain{7} should come out of the local duality machine.
\note{« $F'$ » : exposant frappé biffé et remplacé par un accent, comme p. 63. « SGA~7 » : le chiffre frappé (6 ?) est surchargé à la main ; la lettre a parlé de SGA~7 à ce propos p. 60.}

\add{(d) You may ask Deligne if he did'nt dive into questions (D) and (E) lately.}

\page{67}
\note{pagination de la lettre : « -22- », surchargée en 52.}
\marginal{(App.~14)}
(F) \emph{Significance and limitations of the fppf topology}

Since the attempts of Serre to find a « Weil cohomology » by using the cohomology of a scheme with coefficients not only discrete $\mathbb{Z}/p^n\mathbb{Z}$ $(n \to \infty)$ or $\mu_{p^n}$ $(n \to \infty)$, but also continuous (for example $W_n$, $n \to \infty$), which give good results for recovering a correct $H^1$, during numerous years I have come upon the impression, which I have tried in vain to make precise, that a \add{correct « $p$-adic »} « Weil cohomology », \struck{« $p$-adic »}, in the case $p > 0$ \add{and $k$ of char.~$p$}, should come, in one way or another, from the fppf cohomology, for finite coefficients for example, or more general coefficients, e.g. algebraic groups over $k$. The construction in (B) of the local jacobian complex was, of course, related to this hope: the homology might reveal what is hidden to us in cohomology! For some time now, one has at one\add{'}s disposal the formalism of crystalline cohomology, and one knows (Berthelot) that (at least for $X$ projective and smooth) it has the correct properties. If one uses that as a kind of standard by which to « measure » the other cohomologies, one finds that the part of the crystalline cohomology $H^i_{\mathrm{cris}}(X)$ which could be described in terms of fppf cohomology of $X$ with coefficients in \struck{an} algebraic $k$-group\add{s}. \struck{comes from} \add{is} \add{only} a small part of $H^i$; more precisely, using the very rigid supplementary structure of the $H^i$ (modules of finite type on the ring $W(k)$ of Witt vectors) which comes from the existence of the Frobenius homomorphism \add{(an isogeny)} $H^i \xrightarrow{F} H^i$ (semi-linear), one finds that \struck{they remain} \add{one keeps} always in the part « of slope $\leqslant 1$ » (although the possible slopes vary between $0$ and $i$ …). This explains why for $i = 1$ one can obtain via fppf a correct $H^1$, although for $H^2$ already all the attempts have been unfruitful.
\note{« $\mu_{p^n}$ » : l'indice frappé est surchargé, correction répétée en marge. Le $F$ au-dessus de la flèche de Frobenius est tracé à la main.}

In truth, one conjectures that \emph{all} the part\add{s} \struck{of points} \add{of slope} $\leqslant 1$ in $H^i_{\mathrm{cris}}$ come\struck{s} from fppf. But I have completely lost contact with these questions -- people such as Mazur, Katz, Messing -- and of course Deligne -- should be knowledgeable as to the present states of these questions.
\marginal{espacer}

\marginal{(App.~15)}
Your question 7 seems to indicate that there is a misunderstanding on your part on the significance of the « homotopy type » of $X$, for $X$ a topos

\page{68}
\note{pagination de la lettre : « -23- », surchargée en 53.}
(for example the étale topos of a scheme). Doubtless you must be confusing the homotopical algebra which one can perform on $X$, using semi-simplicial sheaves, \struck{the} stacks of all kinds, the relations between these\add{,} and the other point of view according to which $X$ (with its very rich structure of topos) virtually disappears so as to become no more than a pale element of a « homotopical category » (or pro-homotopical), deduced from the topos by a very thorough process of « localisation »\add{.} \struck{At the} \add{At} first sight, all that still \struck{comes from this} \add{remains with} poor \add{stripped} $X$, are the $\pi_i$ -- and its cohomology groups with constant coefficients -- or at the worst twisted constant coefficients. When one digs more into this definition of « what is left to this poor $X$ » one falls precisely on \emph{the locally constant $n$-stacks} (as an $f : X \to X'$ which is a homotopy equivalence induces a $(n+1)$-equivalence between the categories of locally constant $n$-stacks on $X$ and on $X'$) -- which of course contain the abelian chain complexes of length $n$ of sheaves with locally constant cohomology sheaves, and the hyper-cohomology of these. It is thus that one arrives at this triangle of objects which mutually determine each other:

\begin{tikzcd}[column sep=tiny, row sep=2pt, nodes={font=\scriptsize}]
 & \text{topos (or topological space} & \\
 & \text{or semi-simplicial complex)} & \\
 & \text{\emph{modulo} $n$-homotopy} \arrow[dddl, leftrightarrow] \arrow[dddr, leftrightarrow] & \\
 & & \\
 & & \\
\text{$n$-groupoids} & & \text{``special'' $n$-categories} \\
\text{(up to $n$-equivalence)} \arrow[rr, leftrightarrow] & & \text{(up to $n$-equivalence).}
\end{tikzcd}

\note{chaque sommet du triangle occupe plusieurs lignes dans la frappe ; la disposition est ici reproduite ligne par ligne, les flèches partant des dernières lignes.}

One says that an $n$-category $E_n$ is « special » (or $n$-\emph{galois}) if it is $n$-equivalent to the category of locally constant $(n-1)$-stacks on an appropriate topological space (or a topos), or, what \struck{must} \add{should} be equivalent, if it is $n$-equivalent to the category of $n$-functors $G_n \to (n\text{-}\mathrm{Cat})$, where $G_n$ is an $n$-groupoid. If $X$, $G_n$, $E_n$ correspond in this way, one calls $G_n$ \emph{the fundamental $n$-groupoid} of $X$, or of $E_n$, or says that $E_n$ is \emph{the category of local} \struck{$n$}\add{$(n-1)$}\emph{-systems} on $X$, or on $G_n$, or that $X$ is \emph{the geometric realisation} of $G_n$ or of $E_n$\add{.} In analogy with the familiar case $n = 1$, \struck{$E_n$ can be} \add{it should be possible to} interpret\struck{ed} \add{$G_n$} as the full sub-$n$-category of $\mathit{Hom}_n(E_n, (n\text{-}\mathrm{cat}))$ formed by the $n$-functors $E_n \to (n\text{-}\mathrm{cat})$ satisfying certain exactness properties (one
\note{les soulignements sont ceux de la frappe ; dans « or of $E_n$ » la lettre $E$ est retouchée à la main.}

\page{69}
\note{pagination de la lettre : « -24- », surchargée en 54.}
feels like saying: which commute with finite $\varprojlim$ and arbitrary $\varinjlim$); but this raises the disquieting vision of $n$-limits in $n$-categories. (N.B. The case $n = 2$ begins to become famliliar to us …). It is prudent in all of this to suppose that $X$ is « locally homotopically trivial », which ensures the pro-simplicial set which Artin-Mazur associate to it (with the help of nerves of hyper-coverings) is essentially constant in the ordinary homotopy category -- thus $X$ defines a homotopy type in the usual sense. This is surely \emph{not} the case for the étale topos of a scheme. In such case, the fundamental $n$-groupoid should be conceived as a \emph{pro-$n$-groupoid} (nothing surprising in that, in view of the familiar theory of $\pi_1$), and $E_n$ as an (ind)-$n$-category (the ind-structure will correspond to the exigencies of local triviality for a variable $n$-stack, relative to coverings more and more fine on $X$).
\note{« famliliar » : ainsi dans la frappe.}

\marginal{(App.~16)}
I nevertheless understand your instinctive resistance to conceive this extreme stripping of a beautiful topos $X$, to the point of retaining only the meagre homotopy type. Even more\add{,} I am persuaded that going to the root of this instinctive resistance, one arrives at a generalisation and deepening of the notion of « homotopy type », and to bring new grist to the mill of the development of a good homotopical yoga. Here is what I have in mind.

Let us speak first of sheaves (of sets, or of modules, etc.) instead of stacks, for simplicity, and place ourselves in the étale topos of a scheme. The locally constant sheaves -- modulo a supplementary condition of finiteness which is sufficiently anodyne -- form the easiest of the \emph{constructible} sheaves, for the definition of which they serve as models. Supposing $X$ coherent (= quasi-compact and quasi-separated), then the general constructible sheaves are those for which there exist a finite partition $X = \bigcup_{i \in I} X_i$ of $X$ into « cells » or « strata » $X_i$, each locally

\page{70}
\note{pagination de la lettre : « -25- », surchargée en 55.}
closed and constructible, such that the restriction of $F$ to every $X_i$ is locally constant (also a finiteness condition …). Thus the category of constructible sheaves on $X$ (which gives back the category of all sheaves on passing to a category of ind-objects …) may itself be thought of as an inductive limit of categories associated to finer and finer partitions on $X$. One can then, for such a fixed partition $P$, set out to study the category of sheaves (or complexes of sheaves, or stacks) which are « $P$-constructible » (or, more generally, which are « locally constant » on every $X_i$). These categories will not have truly satisfying structures unless they are stable for the usual operations -- such as $R\,\mathit{Hom}$, or $Rj_*j^*$ where $j : X_i \to X$ is a « cell » of the partition, etc. In fact, if $X$ is excellent and one has resolution of singularities at ones disposal, one knows that the torsion constructible sheaves (under the proviso of being prime to the characteristic) are stable for all these operations -- but not for a finite partition of $X$ fixed once and for all. To have such a finer stability, it is necessary to make some very strict hypotheses of « equi-singularity » on the given stratification of $X$, along the strata. I think nonetheless that a \struck{sufficiency} \add{refinement} of known techniques will show that $X$ admits arbitrarily fine stratifications having these properties of equi-singularity (and with the $X_i$ regular and connected, but this does not matter for our present purpose\struck{s}).

By way of example, suppose that there are \add{just} two \struck{closed} strata, \add{the closed one} $X_0$, and $X_1 = X \setminus X_0$. \struck{After} \add{According to} Artin's devissage, giving oneself a sheaf $F$ on $X$ is equivalent to \struck{being} given a sheaf $F_0 = i_0^*(F)$ on $X_0$, a sheaf $F_1$ $(= i_1^*F)$ on $X_1$, and a homomorphism $F_0 \to i_0^*i_{1*}(F_1)$ \add{$= \varphi(F_1)$}, where $i_0$, $i_1$ are the inclusions $X_0 \xrightarrow{i_0} X \xleftarrow{i_1} X_1$. In order that $F$ should be $P$-constructible, it is necessary and sufficient that $F_0$ and $F_1$ should be locally constant (plus some accessory finiteness conditions …), on $X_0$ and $X_1$ respectively. Then (by virtue of the hypothesis of equi-singularity) the same will be true of \struck{\ill{}} \add{$\varphi(F_1)$,} and the category of sheaves in which we are interested can be expressed entirely in terms of the category of locally constant sheaves on
\note{« equi-singularity » est souligné à la main. Le symbole biffé avant « $\varphi(F_1)$ » (un $X_1$ ?) est noirci. « given » est retouché (« being given » devient « given »). L'alinéa « By way of example » est marqué en marge d'un crochet.}

\page{71}
\note{pagination de la lettre : « -26- », surchargée en 56.}
$X_0$ and $X_1$, i.e. of the mere homotopy type of $X_0$ and $X_1$, except \struck{for the fact} that we must make explicit the nature of the left exact functor $\phi$. I think that this should be possible, in the context of \emph{schemes} in which I am placed (technically rather sophisticated), on introducing an « étale tubular neighbourhood » of $X_0$ in $X_1$ (which is a very \struck{satisfying} \add{interesting} topos, but not associated to a \struck{sheaf} \add{scheme}). But this technical construction is only a paraphrase of an topological intuition extraordinarily simple, which I will make explicit, supposing, to fix the ideas\add{,} that the base field is $\mathbb{C}$ and so one may work with locally compact spaces in the usual sense. The topological idea behind the hypothesis of equi-singularity is that there exists a \emph{tubular neighbourhood} $T$ of $X_0$ in $X$ retracting onto $X_0$ and such that the pair $(X_0, T)$ over $X_0$ should be a locally trivial bundle, i.e. that $T \setminus X_0$ is locally trivial over $X_0$. In fact if $\partial T$ is the « boundary » of $T$, which also should be a locally trivial bundle on $X_0$, then $T$ over $X$ is the conic bundle (= bundle where fibres are cones) $(\simeq (\partial T \times I) \amalg_{\partial T} X_0$ where $I = [0, 1]$, $\partial T \to \partial T \times I$ is defined by $x \mapsto (x, 1)$, and $\partial T \to X_0$ is the projection) then $\mathring{T} = T \setminus X_0 \simeq \partial T \times [0, 1[$ is $X_0$-homotopic to $\partial T$. If $X_0$ and $X_1$ are non singular, then so also will be $\mathring{T}$ and $\partial T$, which are then topologically smooth fibrations on $X_0$. Moreover, putting $\tilde{X}_1 = X_1 \setminus \mathring{T}$, the inclusion $\tilde{X}_1 \to X_1$ is a homotopy equivalence, and $X$ can be recovered, up to homeomorphism, from the diagram of spaces
\note{« topological intuition » est entouré à la main, avec un trait qui le renvoie après « simple » : lire sans doute « an extraordinarily simple topological intuition ». La flèche $\mapsto$ de « $x \mapsto (x,1)$ » est retracée à la main et répétée en marge ; le crochet final de « $[0,1[$ » est corrigé à la main.}

\begin{tikzcd}[column sep=large]
{(\mathring{T} = T \setminus X_0 \simeq)\ \partial T} \arrow[r, "{j\ \text{(inclusion)}}"] \arrow[d, "{\text{fibration}\ p}"'] & {\tilde{X}_1\ (\simeq X_1)} \\
X_0 &
\end{tikzcd}

as an amalgamated sum. In terms of this diagram of spaces, the above functor $\phi$ interprets immediately as
\[
\phi(F_1) \simeq p_*j^*(\tilde{F}_1)
\]
where $F_1 \to \tilde{F}_1$ is the restriction from $X_1$ to $\tilde{X}_1$ (which is an equivalence of categories for the envisaged (locally constant) sheaves). Giving $F = (F_0, F_1, u : F_0 \to \phi F_1)$ can then also be made explicit as giving

\page{72}
\note{pagination de la lettre : « -27- », surchargée en 57.}
\[
F_0 \,,\ \tilde{F}_1 \,,\ \tilde{u} : p^*(F_0) \to j^*(\tilde{F}_1)
\]
where $F_0$ $(\tilde{F}_1)$ are locally constant sheaves on $X_0$ (respectively $X_1$). It is necessary to recall that \struck{up to} here $p$ is a real fibration, and $j$ is an inclusion (in practice, \struck{\ill{}} for the case $X_0$, $X_1$ smooth, \struck{\ill{}} the inclusion of the boundary in a manifold with boundary).
\note{les deux mots biffés de la parenthèse sont noircis ; la virgule après « smooth » est ajoutée à la main.}

If you prefer, one can also take the diagram which is less pretty (but a little more canonical)

\begin{tikzcd}
\mathring{T} \arrow[r, "j'"] \arrow[d, "p'"'] & X_1 \\
X_0 &
\end{tikzcd}

coming essentially to the same thing, as it is formed from spaces homotopic to the preceding one. One can even replace $X_0$ by $T$ ($X_0$ being a deformation retract of it) and write

\begin{tikzcd}
\mathring{T} \arrow[r] \arrow[d, "p''"'] & X_1 \\
T &
\end{tikzcd}

where « literally » $p''$ is now an inclusion, but « morally », it is a \emph{fibration} with \struck{a} very pretty fibres [notably \add{compact,} of finite dimension, and \struck{in fact perhaps} \add{moreover} non-singular varieties -- this is \struck{\ill{}} much better than that which is given by the yoga of Cartan-Serre « every continuous mapping is equivalent to a fibration » …]. This last diagram \struck{is} \add{however has the advantage of being} amenable to a purely algebraic, direct construction, in the context of schemes, once one has developed the construction of étale tubular neighbourhood\add{s} \add{(28)}.
\note{le mot biffé avant « much better » est noirci (« then » ?). Les lignes « construction, in the context … neighbourhoods » sont encadrées à la main, avec une flèche.}

\marginal{espacer}
\marginal{(App.~17)}
The point I wish to come to, is that the consideration of sheaves (or complexes thereof, or $n$-stacks …) which are $P$-constructible on an $X$, where $P$ is a given « equi-singular » stratification, \struck{comes} \add{reduces} in our particular cases to the knowledge of a diagram of ordinary \emph{homotopy types} (or pro-types, if one comes back to the étale topology)
\note{le numéro de l'appel marginal « (App.~17) » est surchargé et lu avec doute ; la parenthèse finale n'est pas suivie de ponctuation.}

\page{73}
\note{pagination de la lettre : « -28- », surchargée en 58.}

\begin{tikzcd}
\Delta \arrow[r, "j"] \arrow[d, "p"'] & X_1 \\
X_0 &
\end{tikzcd}

by taking local coefficients systems (or locally constant $n$-stacks) on the vertices $X_0$, $X_1$, which are related to each other by a homomorphism of compatibility of the type $p^*(F_0) \to j^*(F_1)$. It should be an amusing excercise (which I have not yet done) to verify and to make explicit how the \struck{usual} \add{« six} operations » on sheaves (\add{either} on $X$, or on a \add{subspace which is a} union of strata of $X$) can be expressed in this dictionary, in the case, let us say, of non-singular strata (otherwise, there will be a difficulty with the dualising complexes, which one would prefer to have as objects in our category)\add{,} and \struck{with re-verifying} \add{to reestablish} the known formulae involving these operations. But it appears probable that, to carry out this transcription well, it would be necessary, rather than considering a diagram of type

\begin{tikzcd}
\bullet \arrow[r] \arrow[d] & \bullet \\
\bullet &
\end{tikzcd}

in the homotopical category formed from the category of semi-simplicial sets, to consider the category of diagrams of semi-simplicial sets, and to pass from these to the homotopical category of fractions\add{(29)}.
\note{le sommet du premier diagramme est frappé « $\Delta$ » ; les étiquettes « inclusion » et « fibration » du diagramme de la p. 71 accompagnent $j$ et $p$ dans la frappe. « excercise » : ainsi dans la frappe.}

I have recently more or less made explicit, while thinking on the foundations of « tame topology », (i.e. where one \struck{immediately} eliminates \add{from start} all wild phenomena) how an equisingular stratification, say with non singular strata, of a compact « tame space », gives rise canonically to a diagram of spaces which are manifolds with boundary, the arrows of the diagram being essentially locally trivial fibrations of manifolds with boundary on the others (with fibres \add{which are compact} manifolds with boundary), \emph{and} the inclusion\add{s} of the boundar\struck{y}\add{ies} in these manifolds with boundary [in fact, one finds slightly more general inclusions, certain boundaries which appear being endowed with an « elementary » cellular

\page{74}
\note{pagination de la lettre : « -29- », surchargée en 59.}
decomposition, i.e. the closed strata are again manifolds with boundary which are glued together along common parts of the boundary; and it is also necessary to consider the inclusions of these pieces one in\add{to} another …], and can be reconstituted from this diagram by gluing\add{(30)}. In other words, one has a canonical devissage description of tame compact spaces $X$, eventually endowed with equi-singular stratifications with non-singular strata, in terms of finite diagrams of a precise nature made out of manifolds with boundary. When \struck{in a situation of} \add{we are interested in} sheaves (or complexes of sheaves, or $n$-stacks) which are $P$-constructible on $X$, where $P$ is such a fixed stratification, \struck{known explicitly} \add{these may be described} in terms of the envisaged diagram, of which only the « homotopy type » is to be retained. One \struck{observes} \add{foresees} that the six operations on these sheaves can be translated in an ad hoc manner to this homotopical context. Finally, if instead of having only one compact tame space $X$, one has, let us say, a tame morphism $f : X \to Y$ of such objects, then \struck{by devissage of} \add{by choosing} equi-singular stratifications on $X$ and $Y$ adapted to $f$ (the strata of $X$ being in particular locally trivial fibrations on those of $Y$ …), one should find a « morphism » from the diagram of manifolds with boundary expressing $X$ into that expressing $Y$ (with \struck{some natural} \add{mutual} morphisms which essentially reduce to fibrations of \add{compact} manifolds with boundaries on other such objects) in such a way that the four operations $Rf_*$, $Rf_!$, $Lg^*$, $g^!$ between $P_X$- and $P_Y$- constructible sheaves on $X$ and $Y$ (or on locally closed sub spaces $X'$, $Y'$ which are union of strata, such that $f$ induces $g : X' \to Y'$) can be expressed in terms of standard operations between the mere homotopy type\add{s}. Finally, all these constructions, still partially hypothetical (there is work on the foundations to be done!) should be able to be paraphrased in the framework of excellent schemes, by making use of the machinery of étale tubular neighbourhoods. In one or other case, the « fine homotopy type » of a tame space, respectively of an excellent scheme, is defined by passage to the limit from « $P$-homotopy type » associated to finer and finer equi-singular stratifications $P$ (with non-singular strata).

\page{75}
\note{pagination de la lettre : « -30- », surchargée en 60.}
This « fine » homotopy type would embody the knowledge, not only of sheaves or locally constant $n$-stacks, but (via a passage to the inductive limit) the knowledge of \emph{all of them}. And it would depend\add{,} in a \struck{convenient} \add{suitable} sense, functorially on $X$. In the case of a scheme of finite type on an algebraically closed field $k$ say, the strongest cohomological and homotopical \emph{finiteness theorem} would be expressed precisely in terms of a fine homotopy type, and would say that the ordinary homotopy types which are their constituents are essentially « finite polyhedra » -- and even compact manifolds with boundary -- or in more precise fashion, \struck{the} \add{their} \struck{\ill{}} profinite completions (in the sense of Artin-Mazur) prime to the characteristic $p$ of $k$ \struck{of} \add{are those of} such \struck{spaces.} \add{polyhedra.} One sees clearly how to begin on such a programme in characteristic $0$, but one for\add{e}sees supplementary amusement, or even \struck{of} mystery, in the case $p > 0$, for the varieties which, even birationally, \struck{do} \add{resist} \struck{not lift} \add{being lifted} to characteristic $0$!
\note{« the ordinary homotopy types … finite polyhedra » est souligné à la main. Le mot biffé après « their » est noirci. La marge porte « of/ », renvoi de la correction « are those of ».}

\marginal{(App.~18)}
From these essentially geometric thought\add{s}, I could not at this moment draw up a precise programme for developing \struck{an} adequate algebraic structure\add{s} to express them. I restrict myself to several marginal remarks.

For a long time I have been intrigued by the idea of a « linearisation » of an (ordinary) homotopy type, i.e. questions of the type: if $X$ is a homotopy type, how much cohomological information of the type: cohomology of $X$ with variable coefficients $M$ (constant or twisted constant), multiplicative structure $H^i(X, M) \times H^j(X, N) \to H^{i+j}(X, M \otimes N)$, then eventually other cohomology operations -- is it necessary to have to reconstruct entirely the homotopy type\add{?} (say, in this preliminary pre-derived category approach, \struck{in order to describe} \add{assuming given} the fundamental group $\pi_1$, and therefore the category of constant twisted coefficients (= $\pi_1$-modules), the functors $H^i(-, M)$ over these, together with the structure of cohomological functors relative to exact sequences, the structure of \add{cup-}product, etc. -- related by

\page{76}
\note{pagination de la lettre : « -31- », surchargée en 61.}
certain formal properties\struck{?}?) Once one has at one\add{'}s disposal the language of derived categories: the sub-category of the derived category of abelian complexes\add{ on $X$}, formed from complexes \struck{of} \add{the} sheaves of cohomology \add{of} which are locally constant on $X$, with its triangulated structure and its multiplicative structure $\overset{\mathbb{L}}{\otimes}$ (and eventually $\mathbb{R}\mathrm{Hom}$ …) gives a more satisfying candidate for hoping to recover the homotopy type. I \struck{have little idea} \add{don't really know} if this suffices \struck{to permit} the recovery \struck{without effort} \add{indeed} \add{(31)}, but on the other hand I have no doubt that on \struck{getting to the bottom of the} \add{pursuing} « linearisation » \add{to the end}, that is to say by going to the \emph{non-abelian} framework, and working with the $(n+1)$-category (without any supplementary structure on it!) of locally constant $n$-stacks of constructible sheaves on $X$, for all $n$, one manages to reconstruct the homotopy type via its fundamental $\infty$-groupoid, as explained in my previous letter and recalled in this one. (This signifies in particular that all the possible and imaginable cohomology operations are already included in the data furnished by such a system of $n$-categories …).
\marginal{(31)}
\note{l'appel (31) est répété en marge. Le signe $\infty$ de « $\infty$-groupoid » est retracé à la main.}

Similarly, the more elaborate homotopy type, which are related to certain finite diagrams, which one can associate to certain types of stratification $P$ of tame topological spaces $X$, let's say, should correspond in as perfect a fashion to the $(n+1)$-category of $n$-stacks on $X$ which are locally constant on each of the strata of $P$ (say: which are subordinated to $P$). If the above \struck{machinery of a} \add{description of homotopy types by the} « locally constant derived category » \struck{goes through, one} \add{was valid indeed, one} \add{would} expect\struck{s} to recover here the mixed homotopy type from the \add{corresponding} sub-category of the derived category of abelian sheaves on $X$, provided by the complexes which have locally constant cohomology on each of the strata -- with also the operations $\overset{\mathbb{L}}{\otimes}$, $\mathbb{R}\mathrm{Hom}$, plus in case of need, the four operations $\mathbb{R}g_!$, $\mathbb{R}g_*$, $Lg^*$, $g^!$ for the induced $g : Z' \to Z''$ of the various locally closed unions of strata … . The problem here is that we don't at present even know what is a triangulated category, not \struck{even} \add{any more than} what is it\add{'}s non-commutative version, described probably more simply and more fundamentally\add{:} \struck{as} a « homotopical category » with operations of \add{taking} « fibr\struck{ations}\add{es} » and « cofibr\struck{ations}\add{es} »\add{(32)}.
\note{« one » est ajouté à la main dans la marge droite, à la fin de la ligne biffée.}

\page{77}
\note{pagination de la lettre : « -32- », surchargée en « 61 bis » (lecture incertaine).}
\struck{(N.B. A first approach has been developed by Quillen at the beginning of the years « 70 ».)}
\note{la parenthèse « N.B. … » est biffée à la main de traits obliques.}

It is surely time that I finish this « lettre-fleuve », which is becoming more and more vague. Just one question: what is this marvellous formula of Bloch -- Quillen to which you allude, \struck{and} of which I have never heard, and which makes my mouth water?

Very cordially \struck{to} your\add{s},

Alexandre Grothendieck.
\note{fin de la lettre à L. Breen du 17/19 juillet 1975. Les pp. 2–11 de la lettre, absentes du lot 3, ne figurent pas non plus dans ce lot.}

\page{78}
\note{tapuscrit « NOTES to Chapter I and Appendix » (sa p. 62), notes (1) à (11) aux sections du chapitre I de \emph{Pursuing Stacks} et à l'appendice (lettres à L. Breen), avec corrections et ajouts de sa main (la note (11) est « Added 23.2.83 »). Texte de \emph{Pursuing Stacks}, non transcrit ici : voir G. Maltsiniotis (éd.), \emph{À la poursuite des champs}, vol.~I, SMF, Documents mathématiques 20, 2022. Les notes renvoient aux sections 9, 11, 12, 13 et 18, toutes dans l'intervalle imprimé (1–91).}

\page{79}
\note{suite des notes au chapitre I de \emph{Pursuing Stacks} (sa p. 63), notes (12) à (21), avec ajouts de sa main ; renvois aux sections 13, 18 et 90 (et à « section 11 » p. 78). Non transcrit : voir Maltsiniotis (éd.), \emph{À la poursuite des champs}, vol.~I, SMF 2022.}

\page{80}
\note{suite des notes au chapitre I de \emph{Pursuing Stacks} (sa p. 64), notes (22) à (25) ; la note (22) mentionne « this letter to Larry Breen », la note (24) un ajout de sa main renvoyant à une section dont le numéro, surchargé, se lit 69 ou 59 (« derivator »), la note (25) la section 9. Non transcrit : voir Maltsiniotis (éd.), \emph{À la poursuite des champs}, vol.~I, SMF 2022. La note (25) se poursuit au-delà de ce lot.}

\end{document}
